\documentclass[../DoAn.tex]{subfiles}
\begin{document}

\begin{center}
    \Large{\textbf{ABSTRACT}}\\
\end{center}
\vspace{1cm}
The management of personal finances has become an urgent necessity, especially for young adults – a demographic often characterized by limited income and a lack of established, sound spending habits, leading to the risk of financial mismanagement. Currently, many individuals still rely on manual tools such as notebooks or Excel spreadsheets to track their expenditures. However, these methods expose several serious limitations: they lack the automation needed for seamless recording, present difficulties in data aggregation, and fail to provide the necessary visual analytics that allow users to evaluate and adjust their financial behavior. Stemming from this reality, the final project opts for the approach of developing a Personal Finance Management Application in the form of a web/mobile application. Its goal is to completely replace manual methods and offer immediate, automated data processing capabilities. The solution is built upon a centralized database and a user-friendly interface, focusing on core functionalities: recording daily income and expenses; setting specific budget categories; and generating financial reports through intuitive charts. The system enables users to easily track cash flow, compare actual spending against set budgets, and quickly identify and rectify unreasonable expenditures. The outcome of the project is a complete and stable system that helps users actively manage their personal finances more effectively, contributing to the formation of disciplined spending habits. The main contribution of this thesis is the successful construction of an application model that is simple, user-friendly, yet highly scalable, making it ready for the future integration of advanced features, notably smart spending suggestions based on historical user data.

\end{document}